%% Pitts's alternative contextual first-order presentations, with full proofs.

\subsection{Adjoint presentations of the contextual calculus}
\label{sec:p10-adjoint-first-order-calculus}

\begin{hbdefinition}{def:p10-contextual-adjoint-presentations}{The common basis and two presentations}
Use the language, finite variable contexts, finite hypothesis lists,
and arbitrary theory \(\mathcal T\) of
Definition~\ref{def:p10-contextual-first-order}.
For a binary formula abstraction \(\psi=(x)(x')e\), write
\(\psi(M,N)=e[M/x,N/x']\), with simultaneous capture-avoiding
substitution; \(\psi(x,x)\) denotes its diagonal instance.
The \emph{common basis} consists of the two theory-axiom rules,
Eq-Refl, Eq-Sym, Eq-Trans, Eq-Subst, Id, Cut, Contract, Exchange,
Weaken, and Subst of Sheet~\ref{sheet:contextual-first-order}.
In particular, Exchange exchanges hypotheses, Weaken appends variable
declarations, and the two substitution rules retain their displayed
shared prefix. Contexts remain ordered lists.

For this comparison only, write \(\mathrm{ND}\) for the common basis
together with the propositional, quantifier, and equality-formula
rules of Sheet~\ref{sheet:contextual-first-order}.
Write \(\mathrm{ADJ}\) for the same basis together with all fifteen
directions of Sheet~\ref{sheet:contextual-adjoint-first-order}.
These are local abbreviations for two presentations of derivability.
The notation \(\mathcal T\vdash_{\mathrm{ND}}J\), or
\(\mathcal T\vdash_{\mathrm{ADJ}}J\), means that the judgment \(J\)
has a finite derivation from that same set of theory axioms in the
indicated presentation.
\end{hbdefinition}

Pitts displays seven \emph{bi-rules}: all sequents above a double
line are derivable if and only if the sequent below it is derivable.
This is a statement about derivability, distinct from the formula
constructor \(\Rightarrow\) and from object-language bi-implication.
A bi-rule with two upper sequents gives one downward inference and
two separate upward inferences. The empty upper side of
\(\mathsf{false}\)-Adj gives just one zero-premise inference.
The following sheet displays every such direction of
\cite[Figures~4, 6, and~8, pp.~48, 52, and~55]{Pitts1995Draft}.
The words ``down'' and ``up'' in its rule names refer to the source
double line.

\showcalculus{contextual-adjoint-first-order}

The retained list \(\Phi\) in \(\vee\)-Adj and
\(\exists\)-Adj is part of the displayed source presentation.
Its presence expresses the stability needed when additional
hypotheses are retained. In \(=\)-Adj, \(\Phi\) may depend on
the first equality variable \(x\); only the added variable \(x'\)
is fresh for it. These features agree with the respective semantic
conditions of distributivity, Frobenius, and parameterized equality
in \cite[pp.~48--49, 52--53, 55--56]{Pitts1995Draft}.

\begin{hbtheorem}{thm:p10-contextual-predicate-transport}{proved=annex:p10.adjoint.1}{Syntactic transport of predicate sequents}
Consider any presentation containing the common basis and the
natural-deduction introduction and elimination rules for
\(\Rightarrow\) and \(\forall\).
Let \(\Gamma=(x_1:\sigma_1,\ldots,x_n:\sigma_n)\), let
\(\Phi\vdash\psi\,[\Gamma]\) be derivable from \(\mathcal T\),
and let \(t_i:\sigma_i\,[\Delta]\) be a well-sorted term for
each declaration of \(\Gamma\), in an arbitrary finite context
\(\Delta\). Then
\[
       \Phi[\mathbf t/\mathbf x]\vdash
       \psi[\mathbf t/\mathbf x]\,[\Delta]
\]
is derivable from the same \(\mathcal T\).
Substitution is simultaneous and capture-avoiding.
In particular, consistent variable renaming and permutation of
variable contexts are admissible for predicate sequents.
The statement permits empty \(\Phi\), empty \(\Gamma\), and
arbitrary sets of theory axioms. It asserts no corresponding
transport rule for external term-equation judgments.
\end{hbtheorem}

This syntactic theorem uses implication discharge and universal
closure; it does not treat semantic product isomorphisms as
inference rules or assume that theory axioms are closed under
renaming. The distinction matters because the source contexts and
shared-prefix substitution rules are explicit in
\cite[Definition~2.1.1, p.~6; Figure~1, p.~7;
Remark~3.0.2, p.~14; \S5.1, pp.~40--41]{Pitts1995Draft}.

\begin{hbtheorem}{thm:p10-contextual-adjoint-equivalence}{proved=annex:p10.adjoint.2}{Interderivability of the two contextual presentations}
For every theory \(\mathcal T\) and every well-formed contextual
sequent or external term-equation judgment \(J\),
\[
       \mathcal T\vdash_{\mathrm{ND}}J
       \quad\Longleftrightarrow\quad
       \mathcal T\vdash_{\mathrm{ADJ}}J.
\]
Every direction of each displayed adjoint bi-rule is a derived rule
of \(\mathrm{ND}\), and every logical rule of \(\mathrm{ND}\)
is a derived rule of \(\mathrm{ADJ}\).
Both presentations retain Cut. The correspondence concerns
derivable judgments, without identifying their proof trees.
\end{hbtheorem}

The equivalence asserted in
\cite[pp.~49, 51, and~55]{Pitts1995Draft} is proved here by
explicit finite derivations. Consequently the soundness result
of Theorem~\ref{thm:p10-doctrine-soundness} also applies to
\(\mathrm{ADJ}\), including interpretations with empty sorts.

\begin{hbdefinition}{def:p10-contextual-parameter-free-equality}{The parameter-free equality presentation}
The nearby variant in \cite[Remark~5.6.2, p.~56]{Pitts1995Draft}
restricts the equality bi-rule to an empty retained hypothesis list
and no additional variable context. Denote its two directions
locally by \(\mathrm{Eq0}\), as displayed on the next sheet.
Its binary formula abstraction \(\psi\) has body
\(\psi(x,x')\) in exactly \((x:\sigma,x':\sigma)\), so it
has no additional free parameters. Quantified variables inside
that body are allowed.
\end{hbdefinition}

\showcalculus{contextual-parameter-free-equality}

\begin{hbtheorem}{thm:p10-contextual-equality-reduction}{proved=annex:p10.adjoint.3}{Recovering parameterized equality from the reduced rule}
Over the common basis together with the natural-deduction rules
for implication and universal quantification, the two directions
of full \(=\)-Adj are interderivable with the two directions of
\(\mathrm{Eq0}\).
All other included rules and every theory axiom may be held fixed.
Thus either equality presentation gives the same derivable
judgments under these hypotheses, including when the parameter
context or retained hypothesis list is empty.
No conjunction, disjunction, existential-quantifier rule, or
reflection of formula equality to external term equality is
required in this reduction.
\end{hbtheorem}

The extra implication and universal-quantifier rules in this
theorem are the hypotheses stated in Pitts's remark. They permit
the parameters and hypotheses of the full rule to be represented
inside a formula before the parameter-free rule is applied.
